% POL-2000 · Travail final (A26) · consignes aux étudiant.e.s
% H24 travailFinSession by Adrien Cloutier, kept nearly verbatim. Build: ./build.sh (xelatex, twice)
\documentclass[11pt]{article}
\usepackage{fontspec}
\usepackage[letterpaper,top=0.6in,bottom=0.75in,left=0.8in,right=0.8in,footskip=0.4in]{geometry}
\usepackage{fancyhdr}
\usepackage{lastpage}
\usepackage{xcolor}
\usepackage{graphicx}
\usepackage{tabularx}
\usepackage{array}
\usepackage{enumitem}
\usepackage[most]{tcolorbox}
\usepackage[hidelinks]{hyperref}
\usepackage{parskip}
\usepackage{newunicodechar}
\usepackage{needspace}
\usepackage{ragged2e}
% Espace fine insécable : XeTeX coupait la ligne dessus.
\newunicodechar{ }{\nobreak\hspace{0.16em}\nobreak}

\definecolor{rouge}{HTML}{C6102E}
\definecolor{encre}{HTML}{121314}
\definecolor{papier}{HTML}{F4F4F0}
\definecolor{gris}{HTML}{6B6B6B}

\setmainfont{IBMPlexSans}[Extension=.otf,UprightFont=*-Regular,BoldFont=*-SemiBold,ItalicFont=*-Italic]
\setmonofont{IBMPlexMono}[Extension=.otf,UprightFont=*-Regular,BoldFont=*-SemiBold,ItalicFont=*-Italic]
\newfontfamily\titre{IBMPlexMono}[Extension=.otf,UprightFont=*-SemiBold,BoldFont=*-Bold]

\hypersetup{pdftitle={POL-2000 · Travail final · Consignes},pdfauthor={Laurence-Olivier M. Foisy}}
\color{encre}
\frenchspacing
\setlength{\parskip}{4pt}
\linespread{1.04}

\pagestyle{fancy}
\fancyhf{}
\renewcommand{\headrulewidth}{0pt}
\fancyfoot[L]{\small\ttfamily\color{gris}POL-2000 · Travail final · Automne 2026}
\fancyfoot[R]{\small\ttfamily\color{gris}\thepage{} / \pageref*{LastPage}}

% Titre de section : mono, filet de 2 pt
\newcommand{\etape}[2]{%
  \par\vspace{10pt}\Needspace{4\baselineskip}%
  {\titre\small\color{rouge}#1}\par\vspace{-2pt}%
  {\titre\Large #2}\par\vspace{3pt}%
  {\color{encre}\hrule height 2pt}\vspace{6pt}}

% Encadrés : coins carrés, sans ombre
\newtcolorbox{regle}{enhanced,before upper=\raggedright,sharp corners,boxrule=0pt,leftrule=4pt,colback=papier,colframe=rouge,left=10pt,right=8pt,top=6pt,bottom=6pt}
\newtcolorbox{cadre}{enhanced,sharp corners,boxrule=2pt,colback=white,colframe=encre,left=10pt,right=10pt,top=8pt,bottom=8pt}

\setlist[itemize]{leftmargin=1.3em,topsep=2pt,itemsep=2pt,parsep=0pt}
\setlist[enumerate]{leftmargin=1.6em,topsep=2pt,itemsep=3pt,parsep=0pt}
\definecolor{codefond}{HTML}{EEEEE9}
% Code en ligne : fond gris pâle, comme `code` dans les diapos.
\newcommand{\code}[1]{{\setlength{\fboxsep}{1.6pt}\colorbox{codefond}{\ttfamily\small #1}}}
% Bloc de code : fond gris pâle, cadre encre, filet rouge à gauche (bloc-code des diapos).
\newtcolorbox{codebloc}{enhanced,sharp corners,colback=codefond,colframe=encre,boxrule=1pt,left=10pt,right=8pt,top=4pt,bottom=4pt,before skip=4pt,after skip=6pt,fontupper=\ttfamily\small,overlay={\fill[rouge] (frame.south west) rectangle ([xshift=3.5pt]frame.north west);}}
\newcommand{\bloc}[1]{\begin{codebloc}#1\end{codebloc}}
% Partie du travail : numéro rouge, puis titre
\newcommand{\partie}[2]{\par\vspace{7pt}\Needspace{4\baselineskip}{\titre\large\color{rouge}#1}\hspace{0.6em}{\titre\large #2}\par\vspace{1pt}}

\begin{document}
\raggedright
% Pas de césure : les motifs anglais coupaient mal le français (capa-ble).
\hyphenpenalty=10000
\exhyphenpenalty=10000
\setlength{\parindent}{0pt}

% A26 : version courte, réécrite à partir du travailFinSession H24 d'Adrien Cloutier.
% Mêmes exigences que le plan de cours, deux pages au plus.

% ---------------------------------------------------------------- en-tête
\noindent\begin{minipage}[b]{0.7\textwidth}
{\titre\small\color{gris}POL-2000 · MÉTHODOLOGIE QUANTITATIVE · AUTOMNE 2026}\par\vspace{4pt}
{\titre\fontsize{24}{28}\selectfont Travail final}\par\vspace{2pt}
{\titre\large Consignes}
\end{minipage}\hfill
\begin{minipage}[b]{0.25\textwidth}\raggedleft
\includegraphics[height=0.55in]{ulaval-logo.png}
\end{minipage}
\par\vspace{6pt}{\color{rouge}\hrule height 2pt}\vspace{8pt}

\noindent\begin{tabularx}{\textwidth}{@{}*{4}{>{\raggedright\arraybackslash}X}@{}}
{\titre\small\color{gris}PONDÉRATION} & {\titre\small\color{gris}MODE} & {\titre\small\color{gris}REMISE} & {\titre\small\color{rouge}À REMETTRE}\\
{\titre\Large\color{rouge}25 \%} & Individuel & Format PDF, sur monPortail & \textbf{Vendredi 18 décembre, 23h59}\\
\end{tabularx}

% ---------------------------------------------------------------- description
\vspace{4pt}
Le travail final prolonge votre travail de mi-session. \textbf{L'objectif est de bâtir un article scientifique solide.} Vous travaillez seul.e, avec la base de données de l'Étude électorale canadienne 2025. \textbf{R et Positron sont obligatoires.}

\begin{cadre}
{\titre\small\color{rouge}UN MODÈLE À SUIVRE}\hspace{0.6em}Blais, André et François Gélineau. 2007. « Winning, Losing and Satisfaction with Democracy ». \textit{Political Studies} 55 (2) : 425-441. Vu en séance 8 : exactement ce qu'on vous demande. Lisez-le avant de commencer.
\end{cadre}

% ---------------------------------------------------------------- contenu
\etape{CONTENU}{Ce que votre article doit contenir}

\begin{enumerate}[itemsep=2pt]
\item \textbf{Une question et une hypothèse.} X, Y, une hypothèse et son hypothèse nulle.
\item \textbf{Une revue de la littérature.} Au minimum 8 sources scientifiques, d'où découle votre hypothèse.
\item \textbf{Des données.} Une variable dépendante ordinale, d'intervalle ou de ratio, et au moins un graphique univarié de celle-ci.
\item \textbf{Un graphe orienté acyclique (GOA).} X, Y et au moins trois autres variables qui pourraient biaiser leur relation, dans la base de données ou non.
\item \textbf{Une régression linéaire.} Elle contrôle pour les variables du GOA présentes dans la base de données. Résultats en tableau.
\item \textbf{Une interprétation.} Ce que vos résultats disent de votre hypothèse.
\item \textbf{Des limites.} Les variables du GOA non contrôlées, et une méthode qui estimerait mieux l'effet causal.
\end{enumerate}

\begin{regle}
\textbf{Le reste vous appartient.} Quelle question poser, quelles variables retenir, quels graphiques montrer, comment structurer l'article : ce sont vos choix. Justifiez-les. C'est la qualité de votre raisonnement qui est évaluée.
\end{regle}

% ---------------------------------------------------------------- format
\textbf{Format.} PDF seulement (pas de .docx). 8 pages au maximum, graphiques compris, sans la page couverture ni la bibliographie. Times New Roman 12\,pt, interligne 1,5, références APA (auteur-date), guide de présentation des travaux du Département de science politique. La présentation et la citation des sources comptent dans la note.

\begin{regle}
\textbf{Intelligence artificielle.} Permis, à condition de le déclarer : obtenir de l'aide pour écrire, comprendre ou déboguer du code R. Non permis : l'interprétation des résultats et la rédaction des travaux. Ce que vous affirmez, vous devez pouvoir le défendre. Après la bibliographie, déclarez si vous avez utilisé un outil d'IA. Si oui, lequel, comment et pourquoi ?
\end{regle}

\end{document}
